\chapter*{پیش‌گفتار}
\label{chp_int}
\addcontentsline{toc}{chapter}{پیش‌گفتار}

ساختار نمونهٔ حاضر، مسیر یک پژوهش ریاضی--محاسباتی را از تعریف مسئله تا ارزیابی نتایج نمایش می‌دهد. یک مسئلهٔ مقدار اولیه به‌عنوان مدل معیار در نظر گرفته شده است تا صورت‌بندی ریاضی، تحلیل نظری، تقریب عددی و برآورد پارامتر در چارچوبی یکپارچه ارائه شوند. سادگی مدل معیار امکان مقایسهٔ مستقیم جواب تحلیلی و پاسخ عددی را فراهم می‌کند، بی‌آنکه ساختار فصل‌ها به یک گرایش تخصصی محدود شود.

در فصل نخست، متغیرهای حالت، پارامترها، مشاهدات و تابع هدف تعریف می‌شوند. فصل دوم به خوش‌تعریفی مدل، وابستگی پیوستهٔ جواب به دادهٔ اولیه و شناسایی‌پذیری پارامتر اختصاص دارد. در فصل سوم، همگرایی روش عددی با مطالعهٔ پالایش گام بررسی می‌شود. فصل چهارم یک مطالعهٔ موردی با داده‌های ساختگی را در بر می‌گیرد و میان خطای گسسته‌سازی، خطای برازش و عدم‌قطعیت داده تمایز می‌گذارد. جمع‌بندی، محدودیت‌ها و مسیرهای توسعه در فصل پنجم آمده‌اند.

همهٔ داده‌ها و نتایج عددی این فصل‌ها صرفاً نمونه‌اند و هیچ ادعای تجربی یا پژوهشی مستقلی بر پایهٔ آن‌ها مطرح نمی‌شود. اصطلاحات فارسی نیز با اتکا به واژه‌نامه‌های تخصصی تنظیم شده‌اند \cite{ag85}.
